\documentclass[11pt]{article}
\usepackage[margin=1in]{geometry}
\usepackage{amsmath, amssymb, amsthm}
\usepackage{algorithm}
\usepackage{algpseudocode}
\usepackage{booktabs}
\usepackage{bm}

\title{Methodology: The PRISM Algorithm}
\author{}
\date{}

\begin{document}
\raggedbottom

\maketitle

\section{Problem Formulation}
In Structural Equation Modeling (SEM), the theoretical model is entirely defined by its implied covariance matrix and mean vector. Let $\bm{\theta}$ represent the vector of free parameters. The fundamental equation of SEM states that the population covariance is a function of these parameters: $\bm{\Sigma} = \bm{\Sigma}(\bm{\theta})$. Therefore, preserving the exact sample covariance matrix of a dataset is mathematically identical to preserving all linear structural parameters (e.g., path coefficients, factor loadings, and residual variances) estimated from that data.

Let $\bm{X} \in \mathbb{R}^{n \times p}$ denote an incomplete dataset, and let $\bm{M} \in \{0, 1\}^{n \times p}$ be the missingness indicator matrix ($M_{ij} = 1$ if missing). We define $\bm{X}^{(0)}$ as an initial complete dataset generated via an auxiliary-rich non-parametric imputation method (e.g., Random Forest). While $\bm{X}^{(0)}$ excels at capturing localized, non-linear predictive relationships, single-imputation predictive models inherently attenuate residual variance and violate the global covariance structure implied by the theoretical model.

Let $\bm{\Sigma}^* \in \mathbb{R}^{p \times p}$ be the unbiased target covariance matrix derived from Full Information Maximum Likelihood (FIML) under the Missing At Random (MAR) assumption. The PRISM algorithm (Projecting Imputations on Structural Manifolds) seeks a refined matrix $\bm{X}^*$ that preserves the observed data exactly while deterministically adjusting the imputed values such that the sample covariance of $\bm{X}^*$ matches $\bm{\Sigma}^*$.

\section{Constrained Optimization via Lagrangian Projection}
We frame this structural alignment as a constrained optimization problem. We seek to minimize the squared Frobenius distance between the current sample covariance $\bm{\Sigma}(\bm{X})$ and the target $\bm{\Sigma}^*$, subject to the strict constraint that no observed values are altered. The primal objective function is:
$$ f(\bm{X}) = \frac{1}{2} \big\| \bm{\Sigma}(\bm{X}) - \bm{\Sigma}^* \big\|_F^2 $$

The equality constraint for the observed data is expressed using the Hadamard product ($\odot$):
$$ \bm{X} \odot (\bm{1} - \bm{M}) = \bm{X}^{(0)} \odot (\bm{1} - \bm{M}) $$

Defining the Lagrangian with multiplier matrix $\bm{\Lambda} \in \mathbb{R}^{n \times p}$:
$$ \mathcal{L}(\bm{X}, \bm{\Lambda}) = f(\bm{X}) + \text{Tr}\Big( \bm{\Lambda}^T \big( \bm{X} \odot (\bm{1} - \bm{M}) - \bm{X}^{(0)} \odot (\bm{1} - \bm{M}) \big) \Big) $$

The Karush-Kuhn-Tucker (KKT) stationarity condition requires $\nabla_{\bm{X}} \mathcal{L} = \mathbf{0}$. This condition reveals a critical geometric property: for observed entries ($M_{ij} = 0$), the gradient of the loss is perfectly offset by the dual variables $\bm{\Lambda}$. For missing entries ($M_{ij} = 1$), the dual variable vanishes. Therefore, applying a gradient descent step strictly masked by $\bm{M}$ is mathematically equivalent to projected gradient descent on the feasible affine subspace.

To compute the unconstrained gradient, let $\bm{X}_c = (\bm{I} - \frac{1}{n}\bm{1}\bm{1}^T)\bm{X}$ be the mean-centered data. Applying the chain rule for matrix differentiation, the gradient is:
$$ \nabla_{\bm{X}} f(\bm{X}) = \frac{2}{n-1} \bm{X}_c \big( \bm{\Sigma}(\bm{X}) - \bm{\Sigma}^* \big) $$

\section{Algorithmic Implementation and Guaranteed Convergence}
To ensure monotonic descent and prevent numerical divergence under highly collinear or poorly scaled data, PRISM employs an Armijo backtracking line search. At each iteration, a proposed step size $\alpha$ is evaluated; if the Frobenius loss does not decrease, $\alpha$ is iteratively reduced by a factor of $\tau$ until the sufficient decrease condition is met. Algorithm \ref{alg:prism} details this procedure.

\begin{algorithm}[H]
\caption{The PRISM Algorithm: Masked Covariance Projection}
\label{alg:prism}
\begin{algorithmic}[1]
\Require $\bm{X}^{(0)}$ (Initial imputation), $\bm{M}$ (Missingness mask), $\bm{\Sigma}^*$ (FIML target covariance)
\Require Initial learning rate $\alpha_0$, contraction factor $\tau \in (0,1)$, tolerance $\epsilon$

\State \textbf{PSD Guard:} Ensure $\bm{\Sigma}^*$ is positive semi-definite via spectral decomposition
\State $\bm{X} \leftarrow \bm{X}^{(0)}$

\While{True}
    \State $\bm{X}_c \leftarrow \text{Center}(\bm{X})$
    \State $\bm{\Sigma}_{curr} \leftarrow \frac{1}{n-1} \bm{X}_c^T \bm{X}_c$
    \State $L_{curr} \leftarrow \big\| \bm{\Sigma}_{curr} - \bm{\Sigma}^* \big\|_F^2$
    
    \If{$L_{curr} < \epsilon$} 
        \State \textbf{break} \Comment{Convergence achieved}
    \EndIf
    
    \State $\nabla_{\bm{X}} \leftarrow 2 \bm{X}_c \big(\bm{\Sigma}_{curr} - \bm{\Sigma}^*\big)$
    \State $\alpha \leftarrow \alpha_0$
    
    \While{True} \Comment{Armijo Backtracking Line Search}
        \State $\bm{X}_{prop} \leftarrow \bm{X} - \alpha (\nabla_{\bm{X}} \odot \bm{M})$ \Comment{Masked Lagrangian step}
        \State $\bm{\Sigma}_{prop} \leftarrow \text{Covariance}(\bm{X}_{prop})$
        \State $L_{prop} \leftarrow \big\| \bm{\Sigma}_{prop} - \bm{\Sigma}^* \big\|_F^2$
        
        \If{$L_{prop} < L_{curr}$}
            \State \textbf{break} \Comment{Sufficient decrease met}
        \EndIf
        \State $\alpha \leftarrow \alpha \cdot \tau$
    \EndWhile
    
    \State $\bm{X} \leftarrow \bm{X}_{prop}$
\EndWhile

\Ensure Structure-preserving matrix $\bm{X}^*$
\end{algorithmic}
\end{algorithm}

\section{Extension to Multiple Imputation (PRISM-MI)}
While single imputation via PRISM perfectly recovers point estimates (e.g., structural coefficients), treating imputed values as known quantities inherently downward-biases standard errors. To satisfy Rubin’s Rules (Rubin, 1987) for valid statistical inference, we extend PRISM to Multiple Imputation (PRISM-MI) via asymptotic parameter perturbation.

Let $\hat{\bm{\theta}}$ be the vector of FIML parameter estimates and $\text{ACOV}(\hat{\bm{\theta}})$ be their asymptotic covariance matrix. To generate $m$ proper imputations:
\begin{enumerate}
    \item Draw $m$ perturbed parameter vectors: $\tilde{\bm{\theta}}^{(k)} \sim \mathcal{N}(\hat{\bm{\theta}}, \text{ACOV}(\hat{\bm{\theta}}))$ for $k = 1, \dots, m$.
    \item Reconstruct $m$ distinct model-implied target covariances: $\bm{\Sigma}^{*(k)} = \bm{\Sigma}(\tilde{\bm{\theta}}^{(k)})$.
    \item Execute Algorithm \ref{alg:prism} independently for each $\bm{\Sigma}^{*(k)}$.
\end{enumerate}
This methodology preserves the complex, non-linear relationships captured by the initial machine learning imputation while rigorously enforcing both the structural covariance constraints and the appropriate sampling variability required for inference.

\end{document}
